\documentclass{easychair}

\usepackage{amsmath}
\usepackage{amsfonts}
\usepackage{amssymb}
\usepackage{bussproofs}
\usepackage{stmaryrd}

\theoremstyle{plain}% default
\newtheorem{theorem}{Theorem}
\newtheorem{definition}[theorem]{Definition}
\newtheorem{lemma}[theorem]{Lemma}
\newtheorem{corollary}[theorem]{Corollary}
\newtheorem{problem}[theorem]{Problem}
\newtheorem{example}[theorem]{Example}
\newtheorem{proposition}[theorem]{Proposition}
\newtheorem{remark}[theorem]{Remark}

\newcommand{\rbar}[1]{\mathrel{\bar{#1}}}
\DeclareMathOperator{\dom}{dom}

\title{Higher-order Matching via Intersection Type Inhabitation in Bounded Dimension}

\author{
Andrej Dudenhefner\inst{1}
\and
Aleksy Schubert\inst{2}
}

\institute{
TU Dortmund University, Germany\\
\email{andrej.dudenhefner@cs.tu-dortmund.de}
\and
University of Warsaw, Poland\\
\email{alx@mimuw.edu.pl}
}

\authorrunning{Dudenhefner and Schubert}

\titlerunning{Higher-order Matching via Intersection Type Inhabitation}

\begin{document}

\maketitle

\begin{abstract}
Decidability of higher-order matching was established by Stirling via a sophisticated game-semantic argument.
We propose an alternative specification of higher-order matching as an intersection type inhabitation problem.
Concretely, we use the recently proposed system~\textbf{R}, a syntax-directed intersection type system in which judgments derive vectors of types and binary relations control information flow between coordinates.
Bounding the dimension of type vectors in system \textbf{R} derivations makes inhabitation in system~\textbf{R} decidable.
We outline how a solution to a higher-order matching instance corresponds to an inhabitant in system~\textbf{R} of bounded dimension.
The work is in progress; the central challenge is establishing a dimensional bound that depends only on the matching instance.
\end{abstract}

\section{Introduction}
\label{sect:introduction}

Higher-order matching is the following decision problem: given two simply typed $\lambda$-terms, can the first term be instantiated to be $\beta\eta$-equivalent to the second term?
Higher-order matching was formulated by Huet in the 1970s and constitutes a fundamental decision problem in the simply-typed $\lambda$-calculus, with applications in theorem proving, program transformation, and program synthesis.

Decidability at low orders was established incrementally: order~$2$ by Huet, order~$3$ by Dowek~\cite{dowek1992}, and order~$4$ by Padovani~\cite{padovani2000}.
Decidability of the general problem was settled by Stirling~\cite{stirling2009} using a game-semantic characterization.
Stirling's argument is technically demanding and not easily related to standard type-theoretic constructions.

We propose to pursue a different approach: encoding higher-order matching as an inhabitation problem in a decidable fragment of an intersection type system.
The encoding aims to reduce higher-order matching to inhabitation in a dimensionally bound system~\textbf{R}, a recently introduced intersection type system~\cite{dsr2026sysR}.
This long abstract describes the encoding, identifies the central technical obstacle (a uniform dimensional bound), and discusses possible approaches towards its resolution.

\newpage

\section{System \textbf{R}}
\label{sect:sysR}

We briefly recall system~\textbf{R}~\cite{dsr2026sysR}.
Intersection types $A, B$ and finite sets of intersection types $\sigma, \tau$ are given by
\[
A, B ::= a \mid \sigma \to A
\qquad
\sigma, \tau ::= \{A_1, \ldots, A_n\}, \quad n \geq 0,
\]
where $a$ ranges over type atoms.
A vector $\bar{A} = (A_1, \ldots, A_n)$ has dimension $\dim(\bar{A}) = n$; we write $\bar{A}_{[i]}$ for $A_i$.
The relations $\subseteq$, $\in$, and the operation $\cup$ are extended pointwise to vectors, and denoted $\rbar{\subseteq}, \rbar{\in}$, and $\rbar{\cup}$ respectively.
We set $\bar{\sigma} \Rightarrow \bar{A} := (\sigma_1 \to A_1, \ldots, \sigma_n \to A_n)$.
An \emph{environment}~$\Gamma$ of dimension $d$ assigns to each variable in its domain a vector $\bar{\sigma}$ with $\dim(\bar{\sigma}) = d$.

A binary relation $R \subseteq \{1,\ldots,n\} \times \{1,\ldots,m\}$ acts on a type vector $\bar{A}$ of dimension $n$ by
$R(\bar{A})_{[i]} = \{A_j \mid j\,R\,i\}$,
analogously on vectors of sets, and pointwise on environments.
The operation $(\{A_1\}, \ldots, \{A_n\})\!\!\downarrow\, := (A_1, \ldots, A_n)$ recovers a type vector when each coordinate is a singleton.

\begin{definition}[System \textbf{R} with Dimension Bound $k$]
\label{def:sysR}
~

\centering
{\RightLabel{\textnormal{(Ax)}}
\AxiomC{$\bar{A} \rbar{\in} \bar{\sigma}$}
\AxiomC{$\dim(\Gamma) = \dim(\bar{A}) \leq k$}
\BinaryInfC{$\Gamma, x : \bar{\sigma} \vdash_k x : \bar{A}$}
\DisplayProof}
\quad
{\RightLabel{\textnormal{($\Rightarrow$I)}}
\AxiomC{$\Gamma, x: \bar{\sigma} \vdash_k t : \bar{A}$}
\UnaryInfC{$\Gamma \vdash_k \lambda x.t : \bar{\sigma} \Rightarrow \bar{A}$}
\DisplayProof}~\\[1.5ex]
{\RightLabel{\textnormal{($\Rightarrow$E)}}
\AxiomC{$\Gamma \vdash_{k} t : R(\bar{A}) \Rightarrow \bar{B}$}
\AxiomC{$\Delta \vdash_{k} u : \bar{A}$}
\AxiomC{$R \subseteq \{1, \ldots, \dim(\Delta)\} \times \{1, \ldots, \dim(\Gamma)\}$}
\TrinaryInfC{$\Gamma \cup R(\Delta) \vdash_k t \; u : \bar{B}$}
\DisplayProof}
\end{definition}

The dimension bound $k$ limits the length of all type vectors occurring in a derivation.
The relation $R$ in rule $(\Rightarrow$E$)$ accommodates intersection introduction in a syntax-directed manner: the argument is typed with a vector $\bar{A}$ whose components are combined via $R$.

\begin{example}\label{xmp:selfapp}
With $R := \{(1,1)\}$ in the elimination step we can derive:
\medskip\\
\centerline{\AxiomC{$\{x : (\{\{a\} \to b\})\} \vdash_1 x : R((a)) \Rightarrow (b)$}
\AxiomC{$\{x : (\{a\})\} \vdash_1 x : (a)$}
\RightLabel{\textnormal{($\Rightarrow$E)}}
\BinaryInfC{$\{x : (\{a, \{a\} \to b\})\} \vdash_1 x\,x : (b)$}
\RightLabel{\textnormal{($\Rightarrow$I)}}
\UnaryInfC{$\emptyset \vdash_1 \lambda x.x\,x : (\{a, \{a\} \to b\} \to b)$}
\DisplayProof}
\medskip\\
Even in dimension~$1$, assumptions may carry sets of types of cardinality greater than one.
\end{example}

System~\textbf{R} satisfies subject reduction with respect to $\beta$-reduction in each bounded dimension and corresponds to the standard intersection type system of Barendregt, Coppo, and Dezani-Ciancaglini~\cite{BarendregtCD83} up to top-level subtyping~\cite{dsr2026sysR}.
The inhabitation problem
\begin{quote}
\textsc{Inh}$_{\textbf{R}}$: \emph{Given $\Gamma$, $\bar{A}$, and $k$, does $\Gamma \vdash_k t : \bar{A}$ hold for some term $t$?}
\end{quote}
is the central decision problem of system~\textbf{R}.

\begin{theorem}[Decidability of Inhabitation~\cite{dsr2026sysR}]\label{thm:inh}
\textsc{Inh}$_{\textbf{R}}$ is decidable in 2-EXPTIME.
\end{theorem}

The decision procedure relies on a subformula property: there exists a derivation in $\beta$-normal form in which all coordinates of all type vectors are subformulae of types appearing in the input $\Gamma$ or $\bar{A}$ (or the empty set).
Combined with the dimensional bound $k$, this yields polynomially many distinct vectors per coordinate, which an alternating procedure can traverse traverses in exponential space.

\section{Higher-order Matching via Inhabitation}
\label{sect:homInh}

A higher-order matching instance can be presented as a system of equations of shape either
\begin{equation}\label{eq:hom}
X\,u_i^1\,\ldots\,u_i^n =_{\beta\eta} r_i \qquad (i \in \{1, \ldots, m\})
\end{equation}
or
\begin{equation}\label{neq:hom}
X\,u_i^1\,\ldots\,u_i^n \neq_{\beta\eta} r_i \qquad (i \in \{1, \ldots, m\})
\end{equation}
where the metaviariable $X$ has simple type $\varphi_1 \to \cdots \to \varphi_n \to \varphi$, each $u_i^j$ has simple type $\varphi_j$, each $r_i$ is a simply-typed $\beta\eta$-normal form of type $\varphi$, and free variables are taken from a fixed simply-typed context $\Phi$.
A \emph{solution} is a $\lambda$-term $\lambda x_1 \ldots x_n.t$ of type $\varphi_1 \to \cdots \to \varphi_n \to \varphi$ such that, after substituting it for the metaviariable, every equation and disequation holds.

The encoding combines two ingredients.
\emph{Characteristic typings} capture $\beta\eta$-equality at type~$\varphi$ via intersection-type derivability.
The collection of these typings into a single vector judgment then \emph{combines} into one inhabitation problem.

\subsection{Characteristic Typings}

For each $\beta$-normal $\eta$-long term $r$ with $\Phi \vdash_{\textnormal{STLC}} r : \varphi$, we associate a refining environment $\Gamma$ and intersection type $A$, written $\Gamma \prec \Phi$ and $A \prec \varphi$, such that derivability of a typing for $\Gamma \vdash t : A$ uniquely identifies $r$ up to $\beta$-equivalence among $\eta$-long candidates.

The construction lifts to vectors: a vector $\bar{A}$ characterizes a finite set of right-hand sides $r_1, \ldots, r_m$ by stacking, with $\bar{A}_{[i]} = A_i$ and the corresponding characteristic environments stacked coordinate-wise into $\Gamma$ with $\dim(\Gamma) = m$.

\begin{lemma}[Equality Characterization]\label{lem:char}
For each $\beta$-normal $\eta$-long $r$ with $\Phi \vdash_{\textnormal{STLC}} r : \varphi$, there exist $\Gamma \prec \Phi$ and $A \prec \varphi$ such that:
\begin{enumerate}
\item $\Gamma \vdash r : A$ in system~\textbf{R}
\item for any $\eta$-long $t$ with $\Phi \vdash_{\textnormal{STLC}} t : \varphi$, if $\Gamma \vdash t : A$ then $t \twoheadrightarrow_\beta r$.
\end{enumerate}
\end{lemma}

A complementary disequality characterization is also possible (by exhibiting, for each alternative head shape, an environment ruling it out).

\begin{lemma}[Disequality Characterization]\label{lem:neq_character}
If $r$ is in $\beta$-normal form and $\Phi \vdash_{\textnormal{STLC}} r : \varphi$ is a $\eta$-long derivation, then there exists $\Gamma$, $A$ such that
\begin{enumerate}
\item $\Gamma \prec \Phi$ and $A \prec \varphi$
\item $\Gamma \vdash r : A$
\item for all $t$ such that $\Phi \vdash_{\textnormal{STLC}} t : \varphi$ is an $\eta$-long derivation (not necessarily $\beta$-normal) and $\Gamma \vdash t : A$, then $t \not\twoheadrightarrow_\beta r$
\end{enumerate}
\end{lemma}



\subsection{Reduction to Inhabitation}

For simplicity, let us focus only on $\beta\eta$-equations, since inequations can be dealt with in an analogous manner.

Suppose $s = \lambda x_1 \ldots x_n.t$ is a solution to~\eqref{eq:hom}.
Then $s\,u_i^1\,\ldots\,u_i^n \twoheadrightarrow_\beta r_i$ for every $i$.
By Lemma~\ref{lem:char} together with subject reduction~\cite{dsr2026sysR}, this is equivalent to the existence of some dimension bound $k$ such that
\begin{equation}\label{eq:applied}
\Gamma_{[i]} \vdash_k (\lambda x_1 \ldots x_n.t)\,u_i^1\,\ldots\,u_i^n : \bar{A}_{[i]} \qquad (i \in \{1, \ldots, m\}).
\end{equation}
By inversion of rule $(\Rightarrow$E$)$, there are some relations $R_i^j$ and vectors $\bar{B}_i^j$ with
\[
\Gamma_{[i]} \vdash_k \lambda x_1 \ldots x_n.t : R_i^1(\bar{B}_i^1) \Rightarrow \cdots \Rightarrow R_i^n(\bar{B}_i^n) \Rightarrow \bar{A}_{[i]},
\qquad
\Gamma_{[i]} \vdash_k u_i^j : \bar{B}_i^j.
\]
Stacking the $m$ derivations into vectors of dimension at most $mk$ and combining the relations yields, for some $R^1, \ldots, R^n$ and $\bar{B}^1, \ldots, \bar{B}^n$:
\[
\Gamma, x_1 : \bar{B}^1, \ldots, x_n : \bar{B}^n \vdash_{mk} t : R^1(\bar{B}^1) \Rightarrow \cdots \Rightarrow R^n(\bar{B}^n) \Rightarrow \bar{A}.
\]
The term $t$ is an inhabitant in system~\textbf{R} of dimension $mk$, with environment and target type computable from $r_1, \ldots, r_m$ and the arguments $u_i^j$.

Due to the characterization theorem, the higher-order matching instance~\eqref{eq:hom} has a solution iff the type $R^1(\bar{B}^1) \Rightarrow \cdots \Rightarrow R^n(\bar{B}^n) \Rightarrow \bar{A}$ is inhabited in the environment $\Gamma, x_1 : \bar{B}^1, \ldots, x_n : \bar{B}^n$ for some dimensional bound.

\subsection{The Dimension Bound}

The remaining ingredient is to compute a sufficient dimensional bound based on the given matching instance.
Without such a bound, the inhabitation problem is undecidable.

For fourth-order matching, an explicit bound is plausible, and the following three avenues appear promising.

\begin{description}
\item[Bounded subject expansion.] Starting from the $\beta\eta$-normal $r_i$ and expanding to a candidate solution $(\lambda x_1\ldots x_n.t)\,u_i^1\,\ldots\,u_i^n$, the dimension grows in a controlled fashion governed by the structure of the $r_i$ and the arity profile of the metavariable.
For order $4$, the expansion shape is sufficiently regular that the dimension is bounded by the product of the number of right-hand sides and a quantity depending on the simple types of the arguments.
\item[Padovani-style saturation.] Padovani's decision procedure for fourth-order matching~\cite{padovani2000} bounds the number of necessary type assumptions at intermediate variables.
This induces a bound on the number of applications of intersection introduction in derivations.
Given a bounden number of sufficient subformulae in a type derivation, this may limit the required dimension. 
\item[Krivine-machine states.] Binary relations in system~\textbf{R} bear a resemblance to relational denotational semantics~\cite{Carvalho18}.
The states reachable in a Krivine-machine evaluation of a candidate solution provide a candidate combinatorial bound on the dimension of any associated derivation.
\end{description}

We conjecture that a similar analysis extends to fifth-order matching along the lines suggested by the type structure of fifth-order metavariables, where intersection introduction is required only at applications headed by a main variable.

Alternatively, TODO: it is interesting to render Strilings argument in the language of intersection type inhabitation

\label{sect:bib}
\bibliographystyle{plain}
\bibliography{bibliography}

\end{document}